\documentclass{article}
\usepackage{amsmath}
\usepackage{amssymb}
\usepackage{amsthm}
\usepackage{mathtools}
\usepackage{geometry}
\usepackage{hyperref}

\newcommand{\rk}{\operatorname{rank}}
\newcommand{\GL}{\operatorname{GL}}
\newcommand{\R}{\mathbb{R}}
\newcommand{\Mat}{\mathrm{Mat}}
\newcommand{\Orb}{\operatorname{Orb}}

\theoremstyle{definition}
\newtheorem{theorem}{Theorem}[section]
\newtheorem{proposition}[theorem]{Proposition}
\newtheorem{lemma}[theorem]{Lemma}
\newtheorem{corollary}[theorem]{Corollary}
\newtheorem{definition}[theorem]{Definition}
\newtheorem{remark}[theorem]{Remark}

\title{Task 2: Endpoint Reduction}
\date{}

\begin{document}
\maketitle

\section{Goal}

Reduce Conjecture A to a statement about the endpoint orbit of the end-to-end map $\mu(W)$.

\section{Endpoint group}

Define
\[
K_S:=\{(g_0,g_L)\in \GL(d_0,\R)\times \GL(d_L,\R) : g_0W_0g_L^{-1}=W_0\}.
\]

\begin{definition}[Endpoint action on function space]
The group $K_S$ acts on the function space
\[
\mathcal F:=\Mat_{d_L\times d_0}(\R)
\]
by
\[
(g_0,g_L)\cdot M:=g_LMg_0^{-1}.
\]
\end{definition}

\begin{remark}
The group $K_S$ is the endpoint shadow of the full stabilizer group $G_d^{W_0}$.
\end{remark}

\section{Basic structure}

\begin{lemma}[Projection to the endpoint group]\label{lem:endpoint-projection}
The map
\[
\pi:G_d^{W_0}\to K_S,
\qquad
\pi(g_0,\dots,g_L):=(g_0,g_L),
\]
is a surjective group homomorphism.
\end{lemma}

\begin{proof}
If $g=(g_0,\dots,g_L)\in G_d^{W_0}$, then by definition
\[
g_0W_0g_L^{-1}=W_0,
\]
so $(g_0,g_L)\in K_S$. Hence $\pi$ is well-defined.

It is a group homomorphism because it is just projection onto the first and last coordinates.

To prove surjectivity, let $(g_0,g_L)\in K_S$. Define
\[
\widetilde g:=(g_0,I_{d_1},\dots,I_{d_{L-1}},g_L)\in G_d.
\]
Since $(g_0,g_L)\in K_S$,
\[
g_0W_0g_L^{-1}=W_0,
\]
so $\widetilde g\in G_d^{W_0}$. By construction,
\[
\pi(\widetilde g)=(g_0,g_L).
\]
\end{proof}

\begin{proposition}[Image of a $G_d^{W_0}$-orbit under $\mu$]\label{prop:image-of-orbit}
For every $W\in \Omega$,
\[
\mu\bigl(G_d^{W_0}\cdot W\bigr)=K_S\cdot \mu(W).
\]
\end{proposition}

\begin{proof}
Let $g\in G_d^{W_0}$. By the equivariance formula from Task~1,
\[
\mu(g\cdot W)=g_L\mu(W)g_0^{-1}=\pi(g)\cdot \mu(W),
\]
so
\[
\mu\bigl(G_d^{W_0}\cdot W\bigr)\subseteq K_S\cdot \mu(W).
\]

Conversely, let $(g_0,g_L)\in K_S$. By Lemma~\ref{lem:endpoint-projection}, choose
\[
\widetilde g\in G_d^{W_0}
\qquad\text{with}\qquad
\pi(\widetilde g)=(g_0,g_L).
\]
Then again by equivariance,
\[
\mu(\widetilde g\cdot W)=g_L\mu(W)g_0^{-1}=(g_0,g_L)\cdot \mu(W).
\]
Thus every point of the $K_S$-orbit of $\mu(W)$ is realized as the product of a point in the $G_d^{W_0}$-orbit of $W$, proving the reverse inclusion.
\end{proof}

\section{Endpoint reduction}

\begin{lemma}[Orbit intersection criterion]\label{lem:orbit-intersection-criterion}
Let $W\in \Gamma_{d,S}^r$. Then the following are equivalent:
\begin{enumerate}
\item there exists $g\in G_d^{W_0}$ such that $g\cdot W\in \mathcal C_{d,S}^r$;
\item there exists $(g_0,g_L)\in K_S$ such that
\[
g_L\mu(W)g_0^{-1}=P_SW^\ast;
\]
\item the product $\mu(W)$ lies in the $K_S$-orbit of $P_SW^\ast$.
\end{enumerate}
\end{lemma}

\begin{proof}
\textit{(1) implies (2).}
Assume that there exists $g\in G_d^{W_0}$ such that
\[
g\cdot W\in \mathcal C_{d,S}^r.
\]
Then by definition of the critical locus,
\[
\mu(g\cdot W)=P_SW^\ast.
\]
Writing $\pi(g)=(g_0,g_L)$ and using Task~1,
\[
g_L\mu(W)g_0^{-1}=\mu(g\cdot W)=P_SW^\ast.
\]

\smallskip
\textit{(2) implies (1).}
Assume there exists $(g_0,g_L)\in K_S$ such that
\[
g_L\mu(W)g_0^{-1}=P_SW^\ast.
\]
By Lemma~\ref{lem:endpoint-projection}, choose $\widetilde g\in G_d^{W_0}$ with
\[
\pi(\widetilde g)=(g_0,g_L).
\]
Task~1 proves that $\Gamma_{d,S}^r$ is stable under $G_d^{W_0}$, so
\[
\widetilde g\cdot W\in \Gamma_{d,S}^r.
\]
Moreover,
\[
\mu(\widetilde g\cdot W)=g_L\mu(W)g_0^{-1}=P_SW^\ast.
\]
Hence $\widetilde g\cdot W\in \mathcal C_{d,S}^r$.

\smallskip
\textit{(2) implies (3).}
If
\[
g_L\mu(W)g_0^{-1}=P_SW^\ast,
\]
then by definition of the $K_S$-action the point $P_SW^\ast$ lies in the $K_S$-orbit of $\mu(W)$. Since orbits are equivalence classes for a group action, this is equivalent to
\[
\mu(W)\in K_S\cdot (P_SW^\ast).
\]

\smallskip
\textit{(3) implies (2).}
If
\[
\mu(W)\in K_S\cdot (P_SW^\ast),
\]
then there exists $(h_0,h_L)\in K_S$ such that
\[
\mu(W)=h_L(P_SW^\ast)h_0^{-1}.
\]
Applying the inverse group element gives
\[
h_L^{-1}\mu(W)h_0=P_SW^\ast,
\]
which is exactly (2) with $(g_0,g_L)=(h_0^{-1},h_L^{-1})\in K_S$.
\end{proof}

\begin{corollary}[Endpoint reduction for Conjecture A]\label{cor:conjA-endpoint-reduction}
Conjecture~A is equivalent to the following endpoint normalization statement:
\begin{quote}
For every $W\in \Gamma_{d,S}^r$, the product $\mu(W)$ lies in the $K_S$-orbit of $P_SW^\ast$.
\end{quote}
\end{corollary}

\begin{proof}
Conjecture~A says exactly that every $W\in \Gamma_{d,S}^r$ satisfies condition~(1) of Lemma~\ref{lem:orbit-intersection-criterion}. The endpoint normalization statement is exactly condition~(3) of the same lemma. Therefore the two statements are equivalent.
\end{proof}

\section{Checks and edge cases}

\begin{remark}[Explicit lift when $L=2$]
If
\[
(g_0,g_2)\in K_S,
\]
then the canonical lift
\[
\widetilde g=(g_0,I_{d_1},g_2)\in G_d^{W_0}
\]
satisfies
\[
\mu(\widetilde g\cdot W)=g_2\mu(W)g_0^{-1}.
\]
So in depth two, the reduction to endpoints is completely literal: the middle layer plays no role in the normalization of $\mu(W)$.
\end{remark}

\begin{remark}[Maximal-rank edge case]
If $r=d_L$, then $P_Q=0$, hence
\[
W_0=\Sigma_{XY}P_Q=0.
\]
Therefore
\[
K_S=\GL(d_0,\R)\times \GL(d_L,\R),
\]
and the endpoint reduction becomes the usual left-right orbit problem for matrices of rank $d_L$.
\end{remark}

\section{Deliverables}

This task is complete: the endpoint group is defined, the image of a $G_d^{W_0}$-orbit under $\mu$ is identified with a $K_S$-orbit, and Conjecture~A is reduced to the endpoint normalization problem for $\mu(W)$.

\end{document}
